\section{Рабочий проект}

\subsection{Спецификация классов и компонентов программной системы}

В данном разделе приведено описание основных классов и модулей, реализованных в ходе разработки серверной и клиентской частей программной системы. Для каждого класса указаны назначение, ключевые атрибуты или поля, а также перечень методов с модификаторами доступа и кратким описанием. Для модулей, состоящих из набора функций, приводится описание каждой функции.

Серверный код разделен на пакеты моделей, сериализаторов, представлений, фильтров, валидаторов, общих фрагментов поведения, а также модули проверки прав доступа и маршрутизации. Клиентский код разделен на модули обмена с сервером, многоразовые элементы интерфейса, страницы маршрутов, контексты глобального состояния и пользовательские хуки.

\subsubsection{Класс User}

Класс User описывает модель учетной записи пользователя системы. Он наследуется от стандартного класса AbstractUser платформы Django и расширяет его полем отчества и вычисляемым полем полного имени, объединяющим фамилию, имя и отчество. Класс используется при аутентификации, разграничении прав доступа и поиске пользователей по полному имени или логину. От базового класса наследуются методы хеширования и проверки пароля.

Описание полей класса User приведено в таблице~\ref{tab:rp_user_fields}.

\begin{xltabular}{\textwidth}{|>{\raggedright\arraybackslash}p{40mm}|X|}
\caption{Дополнительные поля класса User\label{tab:rp_user_fields}}\\
\hline
\thead{Поле} & \thead{Назначение} \\
\hline
\endfirsthead
\continuecaption{Продолжение таблицы~\thetable}
\thead{Поле} & \thead{Назначение} \\
\hline
\finishhead
middle\_name & Отчество пользователя \\
\hline
full\_name & Полное имя, вычисляемое базой данных по фамилии, имени и отчеству \\
\hline
\end{xltabular}

\subsubsection{Класс Project}

Класс Project описывает проект как центральную сущность системы. Хранит название, описание, даты начала и завершения, ссылку на создателя. На уровне базы данных задан индекс по названию проекта для ускорения поиска. Через обратное отношение \texttt{members} класс предоставляет доступ к составу участников проекта.

\subsubsection{Класс Task}

Класс Task описывает задачу проекта. Содержит название, описание, даты начала и завершения, порядковый номер задачи в проекте, ссылки на создателя и на проект. Уникальность пары \texttt{(project, number)} запрещает существование двух задач с одинаковым порядковым номером в одном проекте. На поле названия задачи задан индекс. Через обратные отношения \texttt{outgoing\_dependencies} и \texttt{incoming\_dependencies} класс возвращает связи, в которых задача выступает соответственно предшественником и преемником.

\subsubsection{Класс TaskDependency}

Класс TaskDependency хранит зависимость между двумя задачами. Содержит ссылки на задачу-предшественника и задачу-преемника, а также тип связи. Перечень типов связей вынесен в константы класса.

Константы типов связей приведены в таблице~\ref{tab:rp_dep_consts}.

\begin{xltabular}{\textwidth}{|>{\raggedright\arraybackslash}p{40mm}|>{\raggedright\arraybackslash}p{22mm}|X|}
\caption{Константы типов связей класса TaskDependency\label{tab:rp_dep_consts}}\\
\hline
\thead{Константа} & \thead{Значение} & \thead{Описание} \\
\hline
\endfirsthead
\continuecaption{Продолжение таблицы~\thetable}
\thead{Константа} & \thead{Значение} & \thead{Описание} \\
\hline
\finishhead
FINISH\_\allowbreak TO\_\allowbreak START & FS & Преемник начинается после завершения предшественника \\
\hline
START\_\allowbreak TO\_\allowbreak START & SS & Преемник начинается одновременно с предшественником или позже \\
\hline
FINISH\_\allowbreak TO\_\allowbreak FINISH & FF & Преемник завершается одновременно с предшественником или позже \\
\hline
START\_\allowbreak TO\_\allowbreak FINISH & SF & Преемник завершается не раньше начала предшественника \\
\hline
\end{xltabular}

На уровне таблицы базы данных заданы ограничение уникальности пары \texttt{(predecessor, successor)} и проверочное ограничение, запрещающее равенство этих полей.

\subsubsection{Класс ProjectMember}

Класс ProjectMember реализует отношение многие ко многим между пользователями и проектами с указанием роли участника. Содержит ссылки на проект и пользователя, а также поле роли. Перечень ролей вынесен в константы класса.

Константы ролей приведены в таблице~\ref{tab:rp_member_consts}.

\begin{xltabular}{\textwidth}{|>{\raggedright\arraybackslash}p{40mm}|>{\raggedright\arraybackslash}p{22mm}|X|}
\caption{Константы ролей класса ProjectMember\label{tab:rp_member_consts}}\\
\hline
\thead{Константа} & \thead{Значение} & \thead{Описание} \\
\hline
\endfirsthead
\continuecaption{Продолжение таблицы~\thetable}
\thead{Константа} & \thead{Значение} & \thead{Описание} \\
\hline
\finishhead
OWNER & owner & Владелец проекта -- полный доступ \\
\hline
EDITOR & editor & Редактор -- работа с задачами и зависимостями \\
\hline
VIEWER & viewer & Наблюдатель -- доступ только на чтение \\
\hline
\end{xltabular}

Пара \texttt{(project, user)} ограничена уникальностью -- один пользователь не может состоять в одном проекте дважды.

\subsubsection{Класс UserColorPalette}

Класс UserColorPalette хранит индивидуальные настройки оформления интерфейса пользователя. Связан с пользователем отношением один к одному. Содержит поле режима темы \texttt{theme\_mode} (\texttt{system}, \texttt{light}, \texttt{dark}, \texttt{custom}) и набор полей с цветами отдельных элементов: фона приложения, поверхности, шапки, активного пункта меню, акцента, полос задач на диаграмме Ганта, стрелок зависимостей, линий сетки, линии текущего времени, сообщений об ошибках, опасных действий и успешных операций.

\subsubsection{Класс AuthViewSet}

Класс AuthViewSet наследуется от ViewSet библиотеки DRF и предоставляет действия аутентификации пользователя. Описание методов класса приведено в таблице~\ref{tab:rp_auth_viewset}.

\begin{xltabular}{\textwidth}{|>{\raggedright\arraybackslash}p{30mm}|>{\raggedright\arraybackslash}p{18mm}|X|}
\caption{Методы класса AuthViewSet\label{tab:rp_auth_viewset}}\\
\hline
\thead{Метод} & \thead{Доступ} & \thead{Описание} \\
\hline
\endfirsthead
\continuecaption{Продолжение таблицы~\thetable}
\thead{Метод} & \thead{Доступ} & \thead{Описание} \\
\hline
\finishhead
login & public & Проверяет логин и пароль, открывает сеанс работы пользователя, возвращает данные учетной записи \\
\hline
logout & public & Завершает сеанс работы пользователя, возвращает код 204 без тела ответа \\
\hline
csrf & public & Устанавливает cookie csrftoken для последующих изменяющих запросов \\
\hline
\end{xltabular}

\subsubsection{Класс UserViewSet}

Класс UserViewSet наследуется от ModelViewSet библиотеки DRF и реализует операции над учетными записями пользователей. Метод \texttt{get\_queryset} ограничивает выборку текущим пользователем (суперпользователь видит всех). Метод \texttt{get\_serializer\_class} переключает класс сериализатора в зависимости от действия. Метод \texttt{get\_permissions} разрешает создание учетной записи без аутентификации.

Описание методов класса UserViewSet приведено в таблице~\ref{tab:rp_user_viewset}.

\begin{xltabular}{\textwidth}{|>{\raggedright\arraybackslash}p{42mm}|>{\raggedright\arraybackslash}p{18mm}|X|}
\caption{Методы класса UserViewSet\label{tab:rp_user_viewset}}\\
\hline
\thead{Метод} & \thead{Доступ} & \thead{Описание} \\
\hline
\endfirsthead
\continuecaption{Продолжение таблицы~\thetable}
\thead{Метод} & \thead{Доступ} & \thead{Описание} \\
\hline
\finishhead
get\_queryset & public & Возвращает множество доступных пользователю записей \\
\hline
get\_serializer\_class & public & Выбирает класс сериализатора по действию \\
\hline
get\_permissions & public & Определяет требования к авторизации по действию \\
\hline
create & public & Создает учетную запись пользователя \\
\hline
me & public & Возвращает данные текущего пользователя \\
\hline
search & public & Поиск пользователей по логину и полному имени для добавления участников проекта \\
\hline
color\_palette & public & Чтение и частичное обновление цветовой палитры текущего пользователя \\
\hline
\end{xltabular}

\subsubsection{Класс ProjectViewSet}

Класс ProjectViewSet реализует операции над проектами. Метод \texttt{get\_queryset} возвращает только те проекты, в которых текущий пользователь состоит участником. При создании проекта автор автоматически добавляется в участники с ролью владельца в рамках одной транзакции. Изменение и удаление проекта доступны только владельцу.

Описание методов класса ProjectViewSet приведено в таблице~\ref{tab:rp_project_viewset}.

\begin{xltabular}{\textwidth}{|>{\raggedright\arraybackslash}p{42mm}|>{\raggedright\arraybackslash}p{18mm}|X|}
\caption{Методы класса ProjectViewSet\label{tab:rp_project_viewset}}\\
\hline
\thead{Метод} & \thead{Доступ} & \thead{Описание} \\
\hline
\endfirsthead
\continuecaption{Продолжение таблицы~\thetable}
\thead{Метод} & \thead{Доступ} & \thead{Описание} \\
\hline
\finishhead
get\_queryset & public & Возвращает проекты, в которых пользователь состоит участником \\
\hline
perform\_create & public & Создает проект и добавляет автора в участники с ролью владельца в одной транзакции \\
\hline
perform\_update & public & Изменяет проект, требуется роль владельца \\
\hline
perform\_destroy & public & Удаляет проект, требуется роль владельца \\
\hline
\end{xltabular}

\subsubsection{Класс ProjectTaskViewSet}

Класс ProjectTaskViewSet реализует операции над задачами в рамках проекта. Метод \texttt{initial} выполняется до основного действия и проверяет принадлежность пользователя к проекту и его роль. При создании задачи метод \texttt{perform\_create} вычисляет следующий порядковый номер задачи в проекте и записывает его в одной транзакции с самой задачей. Дополнительное действие \texttt{move} изменяет порядок задач в проекте.

Описание методов класса ProjectTaskViewSet приведено в таблице~\ref{tab:rp_task_viewset}.

\begin{xltabular}{\textwidth}{|>{\raggedright\arraybackslash}p{42mm}|>{\raggedright\arraybackslash}p{18mm}|X|}
\caption{Методы класса ProjectTaskViewSet\label{tab:rp_task_viewset}}\\
\hline
\thead{Метод} & \thead{Доступ} & \thead{Описание} \\
\hline
\endfirsthead
\continuecaption{Продолжение таблицы~\thetable}
\thead{Метод} & \thead{Доступ} & \thead{Описание} \\
\hline
\finishhead
get\_queryset & public & Возвращает задачи проекта, упорядоченные по порядковому номеру \\
\hline
initial & public & Проверяет принадлежность пользователя к проекту и его роль до выполнения действия \\
\hline
perform\_create & public & Создает задачу, присваивает порядковый номер в проекте \\
\hline
move & public & Изменяет порядок задач в проекте, перенумеровывает все задачи в транзакции \\
\hline
\end{xltabular}

\subsubsection{Класс TaskViewSet}

Класс TaskViewSet реализует чтение, изменение и удаление задач по глобальному идентификатору. Запросы метода GET без параметра проекта возвращают задачи, созданные текущим пользователем. Запросы PATCH, PUT и DELETE требуют прав редактора или владельца проекта, которому принадлежит задача.

\subsubsection{Класс ProjectDependencyViewSet}

Класс ProjectDependencyViewSet реализует операции над зависимостями между задачами в рамках проекта. Поддерживаются методы GET, POST и DELETE. Метод \texttt{initial} проверяет роль пользователя: для чтения требуется любое участие в проекте, для изменения -- роль редактора или владельца. При создании зависимости в контексте сериализатора передается объект проекта, что используется валидатором для проверки принадлежности обеих задач этому проекту.

\subsubsection{Класс ProjectMemberViewSet}

Класс ProjectMemberViewSet реализует операции над составом участников проекта. Доступ к операциям над участниками имеет только владелец проекта. Изменение и удаление записи об участнике с ролью владельца запрещено. Дополнительное действие \texttt{leave} позволяет участнику самостоятельно покинуть проект.

Описание методов класса ProjectMemberViewSet приведено в таблице~\ref{tab:rp_member_viewset}.

\begin{xltabular}{\textwidth}{|>{\raggedright\arraybackslash}p{42mm}|>{\raggedright\arraybackslash}p{18mm}|X|}
\caption{Методы класса ProjectMemberViewSet\label{tab:rp_member_viewset}}\\
\hline
\thead{Метод} & \thead{Доступ} & \thead{Описание} \\
\hline
\endfirsthead
\continuecaption{Продолжение таблицы~\thetable}
\thead{Метод} & \thead{Доступ} & \thead{Описание} \\
\hline
\finishhead
initial & public & Проверяет аутентификацию и роль пользователя до выполнения действия \\
\hline
get\_queryset & public & Возвращает состав участников проекта, упорядоченный по идентификатору \\
\hline
create & public & Добавляет нового участника в проект с ролью редактора или наблюдателя \\
\hline
perform\_update & public & Изменяет роль участника, запрещает изменение роли владельца \\
\hline
perform\_destroy & public & Удаляет участника, запрещает удаление владельца \\
\hline
leave & public & Позволяет текущему пользователю выйти из проекта \\
\hline
\end{xltabular}

\subsubsection{Класс PasswordStrengthValidator}

Класс PasswordStrengthValidator реализует проверку сложности пароля по требованиям технического задания. Подключается через настройку AUTH\_PASSWORD\_VALIDATORS Django.

Описание методов класса приведено в таблице~\ref{tab:rp_password_validator}.

\begin{xltabular}{\textwidth}{|>{\raggedright\arraybackslash}p{42mm}|>{\raggedright\arraybackslash}p{18mm}|X|}
\caption{Методы класса PasswordStrengthValidator\label{tab:rp_password_validator}}\\
\hline
\thead{Метод} & \thead{Доступ} & \thead{Описание} \\
\hline
\endfirsthead
\continuecaption{Продолжение таблицы~\thetable}
\thead{Метод} & \thead{Доступ} & \thead{Описание} \\
\hline
\finishhead
validate & public & Проверяет, что пароль состоит из допустимых символов и содержит строчную и заглавную латинские буквы, цифру и специальный символ \\
\hline
get\_help\_text & public & Возвращает описание требований к паролю для пользовательского интерфейса \\
\hline
\end{xltabular}

\subsubsection{Модуль проверки зависимостей задач}

Модуль проверки зависимостей задач содержит функции, которые применяются сериализаторами при создании и изменении задач и зависимостей. Функции модуля приведены в таблице~\ref{tab:rp_dep_module}.

\begin{xltabular}{\textwidth}{|>{\raggedright\arraybackslash}p{52mm}|>{\raggedright\arraybackslash}p{18mm}|X|}
\caption{Функции модуля проверки зависимостей задач\label{tab:rp_dep_module}}\\
\hline
\thead{Функция} & \thead{Доступ} & \thead{Описание} \\
\hline
\endfirsthead
\continuecaption{Продолжение таблицы~\thetable}
\thead{Функция} & \thead{Доступ} & \thead{Описание} \\
\hline
\finishhead
is\_pair\_satisfied & public & Проверяет, удовлетворяют ли даты пары задач выбранному типу связи \\
\hline
check\_dependency\_dates & public & Поднимает исключение валидации, если даты задач не соответствуют типу связи \\
\hline
would\_\allowbreak create\_\allowbreak cycle & public & Поиском в ширину проверяет, образует ли новая зависимость цикл в графе зависимостей проекта \\
\hline
check\_\allowbreak task\_\allowbreak dates\_\allowbreak against\_\allowbreak dependencies & public & Проверяет, что новые даты задачи не нарушают входящих и исходящих зависимостей \\
\hline
\end{xltabular}

\subsubsection{Модуль проверки прав доступа}

Модуль проверки прав доступа реализует функции определения роли пользователя в проекте. Используется классами-представлениями для разграничения доступа. Функции модуля приведены в таблице~\ref{tab:rp_permissions_module}.

\begin{xltabular}{\textwidth}{|>{\raggedright\arraybackslash}p{40mm}|>{\raggedright\arraybackslash}p{18mm}|X|}
\caption{Функции модуля проверки прав доступа\label{tab:rp_permissions_module}}\\
\hline
\thead{Функция} & \thead{Доступ} & \thead{Описание} \\
\hline
\endfirsthead
\continuecaption{Продолжение таблицы~\thetable}
\thead{Функция} & \thead{Доступ} & \thead{Описание} \\
\hline
\finishhead
get\_user\_role & public & Возвращает роль пользователя в проекте или None, если пользователь не является участником \\
\hline
is\_member & public & Истинно, если пользователь состоит в проекте \\
\hline
is\_writer & public & Истинно, если пользователь имеет роль владельца или редактора \\
\hline
is\_owner & public & Истинно, если пользователь имеет роль владельца \\
\hline
\end{xltabular}

\subsubsection{Компонент GanttChart}

Компонент GanttChart реализует основное представление календарного плана проекта -- диаграмму Ганта. Отвечает за построение временной шкалы в пределах дат проекта, отрисовку полос задач, стрелок зависимостей, линии текущего времени и штриховки истекшей части полос. Обрабатывает перетаскивание полос задач, перетаскивание левого и правого краев полосы, перетаскивание строки за ручку для изменения порядка задач, режим добавления связей с точками соединения на торцах полос, масштабирование диаграммы в диапазоне 50--400~\% колесом мыши с удержанием Ctrl, кнопками панели инструментов и жестом сведения пальцев. Допустимый диапазон перемещения каждой полосы ограничивается датами проекта и установленными зависимостями.

\subsubsection{Компонент TaskCard}

Компонент TaskCard отображает детальную карточку задачи. Содержит редактируемые поля названия, описания, дат начала и завершения, а также список зависимостей задачи с возможностью добавления и удаления связей. Поля дат ограничены датами проекта и ограничениями из установленных зависимостей.

\subsubsection{Компонент ProjectMembers}

Компонент ProjectMembers реализует управление составом участников проекта. Содержит список текущих участников с их ролями, форму поиска пользователей по логину или ФИО с выводом не более десяти совпадений, элементы изменения роли участника и кнопки удаления. Для текущего пользователя с ролью редактора или наблюдателя предусмотрена кнопка <<Покинуть проект>>.

\subsubsection{Компонент SettingsPanel}

Компонент SettingsPanel реализует страницу настройки цветовой темы. Содержит переключатель режима темы (системный, светлый, темный, пользовательский) и элементы выбора цвета для отдельных компонентов интерфейса в пользовательском режиме. Применение темы выполняется немедленно через изменение CSS-переменных корневого элемента документа без перезагрузки страницы.

\subsubsection{Контексты клиентского приложения}

Глобальное состояние клиентского приложения вынесено в три контекста React. Описание контекстов приведено в таблице~\ref{tab:rp_contexts}.

\begin{xltabular}{\textwidth}{|>{\raggedright\arraybackslash}p{40mm}|X|}
\caption{Контексты клиентского приложения\label{tab:rp_contexts}}\\
\hline
\thead{Контекст} & \thead{Назначение} \\
\hline
\endfirsthead
\continuecaption{Продолжение таблицы~\thetable}
\thead{Контекст} & \thead{Назначение} \\
\hline
\finishhead
AuthContext & Хранит данные текущего пользователя, реализует операции входа, выхода и регистрации, отслеживает состояние загрузки сеанса \\
\hline
ProjectsContext & Хранит список проектов пользователя и идентификатор последнего открытого проекта, реализует операции создания, обновления и удаления проектов \\
\hline
ThemeContext & Хранит режим темы и набор цветов палитры, применяет цвета к корневому элементу документа через CSS-переменные \\
\hline
\end{xltabular}

\subsubsection{Модули обмена с сервером}

Модули обмена с сервером инкапсулируют обращения клиентской части к серверу. Базовый модуль формирует HTTP-запросы с правильными заголовками и cookie, остальные модули реализуют операции над сущностями системы. Описание модулей приведено в таблице~\ref{tab:rp_api_modules}.

\begin{xltabular}{\textwidth}{|>{\raggedright\arraybackslash}p{46mm}|X|}
\caption{Модули обмена с сервером\label{tab:rp_api_modules}}\\
\hline
\thead{Модуль} & \thead{Назначение} \\
\hline
\endfirsthead
\continuecaption{Продолжение таблицы~\thetable}
\thead{Модуль} & \thead{Назначение} \\
\hline
\finishhead
client.js & Базовая функция HTTP-запроса с подстановкой cookie сеанса и заголовка X-CSRFToken \\
\hline
users.js & Операции над учетными записями: регистрация, чтение, обновление, поиск пользователей, чтение и обновление цветовой палитры \\
\hline
projects.js & Операции над проектами: чтение списка с фильтрацией, создание, обновление, удаление \\
\hline
tasks.js & Операции над задачами проекта: чтение списка, создание, обновление, удаление, изменение порядка \\
\hline
dependencies.js & Операции над зависимостями задач: чтение, создание, удаление \\
\hline
members.js & Операции над составом участников проекта: чтение, добавление, изменение роли, удаление, выход из проекта \\
\hline
dates.js & Вспомогательные функции преобразования дат между ISO-строками и календарным представлением \\
\hline
dependencyHelpers.js & Вспомогательные функции определения типа зависимости по выбранным торцам полос задач на диаграмме \\
\hline
roles.js & Константы ролей участников проекта и функции проверки прав на стороне клиента \\
\hline
\end{xltabular}

\subsection{Системное тестирование}

Системное тестирование программной системы проводилось по сценариям, сформированным на основе прецедентов использования из технического задания. Проверялись основные функции системы: регистрация и аутентификация пользователя, управление проектами и составом участников, работа с задачами и зависимостями, интерактивное редактирование на диаграмме Ганта, настройка цветовой темы.

Тестирование выполнялось в операционной системе Windows 11 с использованием браузера Google Chrome. Серверная часть запускалась локально через uvicorn в режиме ASGI, клиентская -- через сервер разработки React. Все сценарии завершены успешно.

\subsubsection{Сводная таблица тест-кейсов}

Сводная информация о выполненных тест-кейсах приведена в таблице~\ref{tab:rp_testcases}.

\begin{xltabular}{\textwidth}{|>{\raggedright\arraybackslash}p{16mm}|X|>{\raggedright\arraybackslash}p{22mm}|}
\caption{Сводная таблица тест-кейсов\label{tab:rp_testcases}}\\
\hline
\thead{ID} & \thead{Название тест-кейса} & \thead{Результат} \\
\hline
\endfirsthead
\continuecaption{Продолжение таблицы~\thetable}
\thead{ID} & \thead{Название тест-кейса} & \thead{Результат} \\
\hline
\finishhead
TC-01 & Регистрация и вход пользователя в систему & Пройден \\
\hline
TC-02 & Создание проекта & Пройден \\
\hline
TC-03 & Добавление задач и просмотр диаграммы Ганта & Пройден \\
\hline
TC-04 & Создание зависимости между задачами & Пройден \\
\hline
TC-05 & Редактирование задачи через карточку задачи & Пройден \\
\hline
TC-06 & Управление составом участников проекта & Пройден \\
\hline
TC-07 & Настройка цветовой темы интерфейса & Пройден \\
\hline
TC-08 & Адаптивное отображение на мобильном устройстве & Пройден \\
\hline
\end{xltabular}

\subsubsection{TC-01. Регистрация и вход пользователя в систему}

\textbf{Предусловие.} Учетная запись с указанным логином не существует, пользователь не аутентифицирован.

\textbf{Шаги:}
\begin{enumerate}
\item Открыть стартовую страницу системы.
\item Нажать кнопку <<Зарегистрироваться>>.
\item Заполнить форму регистрации: логин, адрес электронной почты, пароль, подтверждение пароля.
\item Нажать кнопку <<Зарегистрироваться>> в модальном окне.
\end{enumerate}

\textbf{Ожидаемый результат.} Система проверяет данные, создает учетную запись, открывает сеанс и перенаправляет пользователя на страницу диаграммы Ганта.

\textbf{Фактический результат.} Пройден.

Внешний вид страницы входа и регистрации представлен на рисунке~\ref{fig:rp_welcome}.

\begin{figure}[H]
\centering

\includegraphics[width=\textwidth]{rp_welcome}
\caption{Страница входа и регистрации в систему}
\label{fig:rp_welcome}
\end{figure}

Интерактивная подсказка к паролю в форме регистрации представлена на рисунке~\ref{fig:rp_password_hint}.

\begin{figure}[H]
\centering
\includegraphics[width=\textwidth]{rp_password_hint}
\caption{Интерактивная подсказка к паролю при регистрации}
\label{fig:rp_password_hint}
\end{figure}

\subsubsection{TC-02. Создание проекта}

\textbf{Предусловие.} Пользователь аутентифицирован.

\textbf{Шаги:}
\begin{enumerate}
\item Нажать кнопку создания проекта в шапке приложения.
\item Заполнить форму: название, описание, дата начала, дата завершения.
\item Подтвердить создание.
\end{enumerate}

\textbf{Ожидаемый результат.} Система создает проект, назначает текущему пользователю роль владельца, переключает контекст на созданный проект и переходит к его диаграмме Ганта.

\textbf{Фактический результат.} Пройден.

Внешний вид окна создания проекта представлен на рисунке~\ref{fig:rp_create_project}.

\begin{figure}[H]
\centering
\includegraphics[width=\textwidth]{rp_create_project}
\caption{Модальное окно создания проекта}
\label{fig:rp_create_project}
\end{figure}

\subsubsection{TC-03. Добавление задач и просмотр диаграммы Ганта}

\textbf{Предусловие.} Пользователь аутентифицирован, открыт проект, пользователь имеет роль владельца или редактора.

\textbf{Шаги:}
\begin{enumerate}
\item Нажать кнопку добавления задачи в шапке приложения.
\item Заполнить поля формы и подтвердить создание.
\item Повторить шаги 1--2 для нескольких задач.
\item Открыть страницу диаграммы Ганта.
\end{enumerate}

\textbf{Ожидаемый результат.} Каждая созданная задача получает следующий порядковый номер в проекте и появляется на диаграмме Ганта в виде горизонтальной полосы с названием и длительностью.

\textbf{Фактический результат.} Пройден.

Диаграмма Ганта с задачами проекта представлена на рисунке~\ref{fig:rp_gantt}.

\begin{figure}[H]
\centering
\includegraphics[width=\textwidth]{rp_gantt}
\caption{Страница диаграммы Ганта проекта}
\label{fig:rp_gantt}
\end{figure}

\subsubsection{TC-04. Создание зависимости между задачами}

\textbf{Предусловие.} В проекте существует не менее двух задач, пользователь имеет роль владельца или редактора.

\textbf{Шаги:}
\begin{enumerate}
\item Активировать режим добавления связей кнопкой панели инструментов диаграммы.
\item Нажать на точку соединения на торце задачи-предшественника.
\item Перетащить до точки соединения на торце задачи-преемника.
\end{enumerate}

\textbf{Ожидаемый результат.} Система определяет тип зависимости по выбранным торцам и сохраняет связь. На диаграмме отображается стрелка соответствующего типа.

\textbf{Фактический результат.} Пройден.

Диаграмма с активированным режимом добавления связей и созданной зависимостью представлена на рисунке~\ref{fig:rp_gantt_link}.

\begin{figure}[H]
\centering
\includegraphics[width=\textwidth]{rp_gantt_link}
\caption{Создание зависимости задач перетаскиванием на диаграмме}
\label{fig:rp_gantt_link}
\end{figure}

\subsubsection{TC-05. Редактирование задачи через карточку задачи}

\textbf{Предусловие.} В проекте существует задача с установленной входящей зависимостью, пользователь имеет роль владельца или редактора.

\textbf{Шаги:}
\begin{enumerate}
\item Нажать на полосу задачи на диаграмме Ганта.
\item Изменить название, описание и даты в открывшейся карточке.
\item Подтвердить сохранение.
\end{enumerate}

\textbf{Ожидаемый результат.} Система сохраняет изменения, диаграмма Ганта обновляется. Поля дат в карточке ограничены датами проекта и ограничениями из установленных зависимостей -- задать дату, нарушающую связь, нельзя.

\textbf{Фактический результат.} Пройден.

Карточка задачи представлена на рисунке~\ref{fig:rp_task_card}.

\begin{figure}[H]
\centering
\includegraphics[width=\textwidth]{rp_task_card}
\caption{Карточка задачи}
\label{fig:rp_task_card}
\end{figure}

\subsubsection{TC-06. Управление составом участников проекта}

\textbf{Предусловие.} Пользователь является владельцем проекта.

\textbf{Шаги:}
\begin{enumerate}
\item Перейти на страницу проекта.
\item В разделе участников ввести логин или ФИО добавляемого пользователя.
\item Выбрать пользователя из списка совпадений и назначить ему роль редактора или наблюдателя.
\item Изменить роль одного из существующих участников.
\end{enumerate}

\textbf{Ожидаемый результат.} Новый участник добавлен в проект с указанной ролью, роль существующего участника обновлена.

\textbf{Фактический результат.} Пройден.

Страница проекта с разделом участников представлена на рисунке~\ref{fig:rp_project_page}.

\begin{figure}[H]
\centering
\includegraphics[width=\textwidth]{rp_project_page}
\caption{Страница проекта с управлением составом участников}
\label{fig:rp_project_page}
\end{figure}

\subsubsection{TC-07. Настройка цветовой темы интерфейса}

\textbf{Предусловие.} Пользователь аутентифицирован.

\textbf{Шаги:}
\begin{enumerate}
\item Перейти в раздел <<Настройки>> через боковую панель навигации.
\item Выбрать пользовательский режим темы.
\item Изменить цвета отдельных элементов интерфейса.
\end{enumerate}

\textbf{Ожидаемый результат.} Выбранные цвета применяются к интерфейсу немедленно без перезагрузки страницы и сохраняются в учетной записи пользователя.

\textbf{Фактический результат.} Пройден.

Страница настроек цветовой темы представлена на рисунке~\ref{fig:rp_settings_page}.

\begin{figure}[H]
\centering
\includegraphics[width=\textwidth]{rp_settings_page}
\caption{Страница настроек цветовой темы}
\label{fig:rp_settings_page}
\end{figure}

\subsubsection{TC-08. Адаптивное отображение на мобильном устройстве}

\textbf{Предусловие.} Пользователь аутентифицирован, ширина окна браузера соответствует размерам экрана смартфона.

\textbf{Шаги:}
\begin{enumerate}
\item Открыть страницу диаграммы Ганта на узком экране.
\item Перейти на страницу проекта и в раздел настроек.
\end{enumerate}

\textbf{Ожидаемый результат.} Раскладка перестраивается под ширину экрана. Горизонтальная прокрутка вне диаграммы Ганта отсутствует, функциональность сохраняется.

\textbf{Фактический результат.} Пройден.

Вид интерфейса системы на мобильном устройстве представлен на рисунке~\ref{fig:rp_mobile}.

\begin{figure}[H]
\centering
\includegraphics[width=0.5\textwidth]{rp_mobile}
\caption{Адаптивное отображение интерфейса на мобильном устройстве}
\label{fig:rp_mobile}
\end{figure}
